(1) 令原胞第 \(n\) 个质量 \(m_1\) 的位移为 \(u_n\)，质量 \(m_2\) 为 \(v_n\)，则 \(m_1\ddot u_n=K(v_n+v_{n-1}-2u_n)\)、\(m_2\ddot v_n=K(u_n+u_{n+1}-2v_n)\)。每个原子受左右两弹簧作用，故对角恢复力系数为 \(2K\)。
(2) 代入 \(u_n=Ue^{i(nka-\omega t)}\)、\(v_n=Ve^{i(nka-\omega t)}\)，行列式为 \((2K-m_1\omega^2)(2K-m_2\omega^2)-4K^2\cos^2(ka/2)=0\)。解得 \(\omega_\pm^2=K(1/m_1+1/m_2)\pm K\sqrt{(1/m_1+1/m_2)^2-4\sin^2(ka/2)/(m_1m_2)}\)。
(3) 在 \(k=0\)，低支 \(\omega_-(0)=0\)、高支 \(\omega_+^2(0)=2K(1/m_1+1/m_2)\)。在区边界 \(k=\pi/a\)，两频率平方分别为 \(2K/\max(m_1,m_2)\) 与 \(2K/\min(m_1,m_2)\)；若质量不同则存在分隔两支的有限带隙，若相同则边界简并。
(4) \(|ka|\ll1\) 时，将根号展开，得 \(\omega_-^2\simeq K a^2k^2/[2(m_1+m_2)]\)，因此声速 \(c=a\sqrt{K/[2(m_1+m_2)]}\)。在 \(m_1=m_2=m\) 时 \(c=(a/2)\sqrt{K/m}\)，恰等于相邻间距 \(a/2\) 的单原子链长波声速；光学分支是采用加倍原胞后布里渊区折叠的结果。
